% Generated draft: requires mathematical and source review.
\documentclass[11pt]{article}
\usepackage[margin=25mm]{geometry}
\usepackage{fontspec}
\setmainfont{Latin Modern Roman}
\usepackage{amsmath,amssymb,mathtools,unicode-math}
\setmathfont{Latin Modern Math}
\usepackage{enumitem,needspace,xurl,hyperref}
\setlength{\parindent}{0pt}
\setlength{\parskip}{5pt}
\setlength{\emergencystretch}{3em}
\newcommand{\N}{\mathbb{N}}
\newcommand{\Z}{\mathbb{Z}}
\newcommand{\Q}{\mathbb{Q}}
\newcommand{\R}{\mathbb{R}}
\newcommand{\C}{\mathbb{C}}
\newcommand{\F}{\mathbb{F}}
\newcommand{\A}{\mathbb{A}}
\newcommand{\PP}{\mathbb{P}}
\newcommand{\E}{\mathbb{E}}
\newcommand{\eps}{\varepsilon}
\newcommand{\dd}{\,\mathrm{d}}
\begin{document}
\section*{UnsolvedMath ID: 30003983}
\textbf{Source problem code:} catalogue entry
\subsection{Definitions and mathematical statement}
The supplied record asks for a characterization of decompositions of an equiangular tight frame into equiangular tight frames for their respective spans. It does not specify the scalar field, the normalization convention, or the meaning of ``nontrivial.'' The following gives a precise finite, unit-norm formulation over each of the fields \(\mathbb K\in\{\R,\C\}\), without asserting that these conventions resolve the source's omissions.

Let \(H\) be a \(d\)-dimensional inner-product space over \(\mathbb K\), where \(1\le d\le N\), and let \(\Phi=(\varphi_i)_{i\in I}\) be an indexed family spanning \(H\), with \(|I|=N\) and \(\|\varphi_i\|=1\) for every \(i\in I\). It is an equiangular tight frame if there exists \(\alpha\in[0,1]\) such that
\[
 |\langle\varphi_i,\varphi_k\rangle|=\alpha
 \qquad(i,k\in I,\ i\ne k),
\]
and
\[
 \sum_{i\in I}|\langle x,\varphi_i\rangle|^2
 =\frac{N}{d}\|x\|^2
 \qquad(x\in H).
\]
Here equiangularity is vacuous for a family with one vector. Whether repeated vectors or repeated one-dimensional subspaces are permitted is not specified in the source.

A decomposition is a partition
\[
 I=I_1\mathbin{\dot\cup}\cdots\mathbin{\dot\cup}I_r,
 \qquad 2\le r\le N,\qquad I_j\ne\varnothing.
\]
For each \(j\), set
\[
 \Phi_j=(\varphi_i)_{i\in I_j},\qquad
 H_j=\operatorname{span}_{\mathbb K}\{\varphi_i:i\in I_j\},\qquad
 d_j=\dim H_j.
\]
The required block condition is that \(\Phi_j\) be an equiangular tight frame for \(H_j\). Under the unit-norm convention above, equiangularity is inherited from \(\Phi\), so the additional condition is
\[
 \sum_{i\in I_j}|\langle x,\varphi_i\rangle|^2
 =\frac{|I_j|}{d_j}\|x\|^2
 \qquad(x\in H_j),\qquad 1\le j\le r.
\]
Disjointness here concerns the index sets; it does not require the subspaces \(H_j\) to be disjoint, orthogonal, or proper subspaces of \(H\).

The intended problem is to characterize those equiangular tight frames admitting such a decomposition that is nontrivial in the source's intended sense. That sense remains unresolved: requiring only \(r\ge2\) permits the partition into singletons whenever \(N\ge2\). The supplied record does not state whether singleton blocks, the all-singleton partition, or any other decompositions are to be excluded. Thus the displayed conditions specify the decomposition relation, but an exact nontrivial-existence predicate requires an additional source clarification. The record also does not explicitly specify finiteness or finite dimensionality; the formulation above is conditional on those standard finite-frame conventions.

The supplied dated curation describes the two-block case as structurally characterized, with the block spans either orthogonal or both equal to the full span, and says that the equiangular special case is correspondingly restricted while the broader decomposition question remains open. These are supplied literature-assessment claims, not independently verified theorems here; their precise hypotheses and the stated equiangular restrictions are not included in the record.
\subsection{Short English statement}
Determine when an equiangular tight frame can be partitioned into smaller families that are themselves equiangular tight frames for their own spans. A precise finite, unit-norm decomposition condition is given above for real and complex spaces separately. The source does not define ``nontrivial''; allowing all singleton blocks makes existence automatic for frames with at least two vectors. The supplied status is \emph{partially solved}, subject to its dated, unverified literature assessment.
\subsection{Sources}
\begin{itemize}
\item Supplied local UnsolvedMath registry record, permanent ID \(30003983\), title \emph{Decomposition of Equiangular Tight Frames}, source problem code \texttt{catalogue entry}; collection: \emph{Oberwolfach Reports Open Problems}. No individual proposer is supplied. The proposed year is recorded as 2018.
\item Supplied source citation: \emph{Mini-Workshop: Algebraic, Geometric, and Combinatorial Methods in Frame Theory (2019)}, Oberwolfach Reports, pp.\ 2769--2803, DOI \texttt{10.4171/owr/2018/46}. Supplied source and background link: \url{https://doi.org/10.4171/owr/2018/46}. The citation's 2019 designation and the report's 2018 designation are retained as supplied.
\item Supplied literature reference: \emph{Oberwolfach Report 46/2018}, \url{https://doi.org/10.4171/owr/2018/46}. Literature checked at: \texttt{2026-08-21}. This date and the assessment are supplied dated curation, not independently verified. The curation labels the problem partially solved, describes the two-block decomposition as structurally characterized with spans either orthogonal or both full, says the equiangular special case is correspondingly restricted, and describes the broader question as remaining open. It explicitly requests independent verification of time-sensitive status and cited claims.
\item The registry assigns statement status \texttt{corrected\_verified} and reports that the official Oberwolfach report was checked directly, with definitions restored and source-dependent wording removed. These are supplied verification assertions; no independent source consultation or mathematical verification is claimed in this draft.
\end{itemize}
\end{document}
